Put $\mathfrak q=(x_1,\ldots,x_d)$ and $R=k[[t_1,\ldots,t_d]]$, with $\mathfrak n=(t_1,\ldots,t_d)$. Since $\mathfrak q$ is $\mathfrak m$-primary, the $\mathfrak q$-adic and $\mathfrak m$-adic topologies agree. Hence $A$ is complete and Hausdorff for the $\mathfrak q$-adic topology, and substitution defines $R\to A$.

Suppose a nonzero series $F$ maps to zero, and let $F_r$ be its first nonzero homogeneous part. The remaining terms have image in $\mathfrak q^{r+1}$, so $F_r(x_1,\ldots,x_d)\in\mathfrak q^{r+1}$. Proposition 11.20 implies that every coefficient of $F_r$ lies in $\mathfrak m$. These coefficients belong to $k$, and $k\cap\mathfrak m=0$, a contradiction. Thus the map is injective.

The quotient $A/\mathfrak q$ has finite length, hence is finite-dimensional over $k$. The graded module $\operatorname{gr}_{\mathfrak q}A$ is generated over $\operatorname{gr}_{\mathfrak n}R=k[t_1,\ldots,t_d]$ by a $k$-basis of $A/\mathfrak q$: every graded piece is generated by monomials in the initial forms of the $x_i$ with coefficients in $A/\mathfrak q$. Since $R$ is complete and $A$ is Hausdorff, Proposition 10.24 applied to the filtration $\mathfrak n^rA=\mathfrak q^r$ makes $A$ a finite $R$-module.
